\section{Comparison of the three supplied proposals}\label{sec:comparison}

\subsection{Snapshot and literal comparison}
The comparison is pinned to ProveIt commit
\begin{center}\small
\hash{da4ab9dd68ccfc47ac852118601eaf1fdae53795}.
\end{center}
Its \file{docs/incoming/} directory contained the three U211 archives and
\file{pa_unary_power_cover.zip} used here~\cite{snapshot}.
All four original archives are retained unchanged in \file{inputs/}.
Their SHA-256 values, original paths, and sizes are recorded in
\file{provenance.json}. The table/source snapshots used by the new construction
are checked directly against the corresponding ZIP members.

\begin{table}[ht]
\centering\small
\begin{tabularx}{\linewidth}{@{}lrrX@{}}
\toprule
Proposal & Rotations & Deleted label & Principal mechanism\\
\midrule
Contextual & 6 & 180 &
Predecessor-tag fold into 163 through helper 211;
also repairs the backward entry \(171/0\) into state 7.\\
Research & 9 & 181 &
Predecessor-tag fold into 220 through helper 129;
retains the earlier helper 211.\\
Rotated & 11 & 218 &
Makes every entry to 218 write a zero immediately left of the head,
then replaces 218 by the existing state 157.\\
\bottomrule
\end{tabularx}
\caption{Distinct 211-state constructions from the same U212 parent.}
\end{table}

All three contain the same ordered 3,909 primitive operations and operands.
Both source-control resolvers produce the same complete ordinal control-flow
graph. However, their native tables differ, and canonical breadth-first
numbering from state 0 shows that their rooted transition graphs are pairwise
nonisomorphic when read symbols, write symbols, and directions are preserved.
These are different machines, not three filenames for one table.

\begin{table}[ht]
\centering\small
\begin{tabular}{lrr}
\toprule
Compared pair & Different common native rows & Different source addresses\\
\midrule
Contextual / Research & 28 & 1,275\\
Contextual / Rotated & 31 & 3,096\\
Research / Rotated & 29 & 3,096\\
\bottomrule
\end{tabular}
\caption{The row counts exclude the two outgoing rows of each differently
deleted label. Full inventories are in \texttt{comparison/census.json}.}
\end{table}

The common parent has native SHA-256
\begin{center}\small
\hash{386ce16259d4c97e4af31b3ba494feaff11a0a1a066e0235dddc154b4cea38c4}.
\end{center}
The three U211 native identities are listed below in abbreviated form;
the full hashes are retained in the comparison files.
\[
 \begin{array}{ll}
 \text{Contextual:}&\texttt{ac75e2fa934e89b3\ldots}\\
 \text{Research:}&\texttt{7edae566b4c2e5c3\ldots}\\
 \text{Rotated:}&\texttt{72234980bd604524\ldots}.
 \end{array}
\]

\subsection{Verification quality and the selected parent}
Each original package was copied to an isolated working directory and its
finite replay command was rerun successfully: Contextual's default reproduction,
Research's \texttt{--quick} reproduction, and Rotated's default verifier.
Those fresh receipts and logs are included. The original full symbolic runtime
suites were not rerun for all three inputs; their packaged full-runtime
receipts remain prior evidence. In contrast, the complete final-table runtime
derivations for both new candidates were freshly regenerated.

The Research table is the direct parent of the new work because its tagged
fold can be preserved while the 218 predecessor condition is rearranged.
It also gives a particularly transparent 210-state intermediate without the
171-entry repair. The Contextual package contributes an independently useful
retagging and repair argument; Rotated contributes the exact 218-to-157
two-cell identity. None of these merits implies that its entire patch can be
copied into another source layout without rechecking the entry conditions.

The common consistency calibration is not a reason to treat all three as
already covered by K215's separate formal endpoint. The linked
K215/MM0 project proves a statement about a different literal table and
ZFC~\cite{k215pa}. Its result motivates the architecture but does not certify
the newly supplied PA machines.

\section{Theory U and the arithmetic program actually preserved}\label{sec:theory}

\subsection{The unary schema}
Theory U uses classical, equality-free first-order logic on nonempty domains,
with one binary relation \(x\in a\). Put
\[
 \Sub(x,a)\;:\!\Longleftrightarrow\;\forall z\,(z\in x\to z\in a).
\]
For every formula \(\varphi(x)\) with no other free variable, its sole schema is
\[
 \forall a\,\exists b\,
 \left[
   \forall x\,(x\in b\to\varphi(x))
   \ \land\
   \forall x\,(\Sub(x,a)\to(\varphi(x)\to x\in b))
 \right].
 \tag{U}
\]
Variables are freshened before insertion. The source \(a\) is not a permitted
free matrix parameter. This is a restriction on free roles, not a restriction
to formulas containing only one variable symbol.

The schema supplies a weak cover: all members of \(b\) satisfy the matrix, and
every subset of \(a\) satisfying it belongs to \(b\). It does not require every
member of \(b\) to be a subset of \(a\), and it does not postulate ambient
extensionality. Replacing the matrix by
\(\Sub(x,a)\land\varphi(x)\) would introduce the forbidden free source \(a\);
that apparent shortcut is not available in U.

\subsection{Why the retained U/PA proof is substantive}
The supplied unary manuscript~\cite{unary} contains a detailed ordinary
interpretation argument. Its upper direction interprets membership as binary
bit membership in PA. Bit inclusion implies \(x\le a\); for each fixed
arithmetical matrix, PA's induction constructs the finite characteristic code
on the interval through \(a\). This is formula-specific finite comprehension,
not a uniform truth predicate, and works inside arbitrary PA models.

For the lower direction, equality is interpreted by coextensionality.
The numerical representatives explicitly include saturation and suitable
first-order order/endpoint conditions. Local predecessor facts establish
induction before total successor existence is used. Comparison and least
failing numerical tuples then allow parameters to be bound inside a legal
unary matrix. Bounded numerical comprehension follows, with its support and
invariance premises stated explicitly. This provides successors, numerical
powers, pair and graph codes, and canonical finite histories for addition and
multiplication. The order of these steps avoids a circular appeal to total
successors.

The manuscript also describes uniform syntactic proof templates: freshening,
substitution, relativization, replacement of terms by functional graphs, and
concatenation of finite derivations. These are primitive recursive syntax
operations, with correctness justified by induction in PA. An independent
review here found no concrete defect in that dependency chain. The review
does not claim to have implemented a complete arbitrary-proof translator or
checked an exported endpoint in a proof assistant.

\subsection{The machine-facing calculus is a different presentation}
The preserved source program searches the inherited numerical calculus
\(\DU\). It uses truth code \(0\), falsity code \(1\), native Cantor-coded
constructors, seven logical rows, a dominated propositional row, modus ponens,
generalization, and a gated cover/vacuity bundle.
Its machine-facing theory allows additional matrix parameters and uses global
irreflexivity. It is calibrated by finite-refutation translations to the
unary U presentation, rather than having literally the same axiom strings.

This distinction is essential. The separate positive dyadic emitter in the
unary-theory ZIP is not the program implemented by these U211, U210, or U209
tables. Its operation counts cannot be used as a runtime or state-count
certificate for them.

The retained role-emitter proof handles small-code aliases, the truth/falsity
codes, the gate, and the possible coincidence of output and source labels.
In the coincident-label case, irreflexivity changes the relevant guard to
emptiness; an empty source supplies the required cover. Padding and fresh
labels supply completeness for arbitrary desired instances.
The new reductions leave this numerical calculus and every ordered source
primitive unchanged.

\section{Source equivalence and contextual decoding}\label{sec:method}
The source is a binary tree whose leaves carry \(\mathrm{CI}\) (increment),
\(\mathrm{CD}\) (guarded decrement), \(\mathrm{CB}\) (forward branch),
\(\mathrm{CC}\) (backward branch), or \(\mathrm{CH}\) (halt).
There are respectively 1,558, 1,175, 80, 1,095, and one such occurrences.
The leaf address encodes control context, so retaining the printed operation
sequence alone is insufficient.

For ordinal \(i\), let \(G(i)\) record the primitive, its operand, and its ordered
successor ordinals. A decrement has zero and positive successors; halt has
none. The first resolver computes successors by binary increment and suffix
cutting, followed by zero-padded leaf decoding. The second resolver uses the
ordinal intervals of enclosing subtrees. Both check all 5,083 nonhalt
successor edges. The arithmetic resolver additionally rejects stray one bits
after the decoded target leaf.

A factor is closed and nonhalting when its ancestor-relative branches are
resolved within the factor and it contains no halt. Closure is checked in
two ways: recursive boundary signatures and independent leafwise counts of
the relevant enclosing zero/one edges. Association of closed sequential
factors preserves the primitive word and resolved ordinal control graph.

\begin{lemma}[Source conjugacy]\label{lem:source}
If two complete source trees have the same ordered primitive word and the
same \(G(i)\) for every ordinal, then from a common ordinal and numerical
register store they have identical stores and ordinals at every source step.
They halt after the same number of source steps.
\end{lemma}
\begin{proof}
Induct on the source-step count. Increment changes the same register in both
programs. A decrement tests the same value, so both programs choose the same
branch and make the same register change. The two branch operations choose
the checked common successor; halt occurs in the same ordinal. No bound on
the register values is used.
\end{proof}

A pure physical decoder has rows \((0,R,a)\) and \((1,R,b)\).
A tagged physical decoder may implement different logical source subtrees
depending on a bit left by its incoming transition. Therefore a source rewrite
must preserve the relevant entry context. This is the key to the new result:
one may alter a branch of a shared decoder that a runtime call cannot execute,
while preserving all branches used by that runtime call.

The source theorem is separate from the final literal-machine theorem.
Some construction intermediates have seven consecutive zero address bits,
which collide with the inherited origin detector. Their source control graphs
are still equivalent, but those intermediate tables are not claimed to
implement a complete correct runtime. The final sources have maximum zero
run six and pass the complete runtime verification directly.
